% =====================================================================
%  FICHE 8 — THÈME 4 — NIVEAU MOYEN — ÉNONCÉS
% =====================================================================
\documentclass[11pt,a4paper]{article}
\usepackage[utf8]{inputenc}
\usepackage[T1]{fontenc}
\usepackage{lmodern}
\usepackage[french]{babel}
\usepackage{amsmath,amssymb,mathtools}
\usepackage{icomma}
\usepackage{siunitx}
\sisetup{output-decimal-marker={,},per-mode=power,group-separator={\,},group-minimum-digits=5}
\usepackage[version=4]{mhchem}
\usepackage{xcolor}
\usepackage{tikz}
\usepackage[most]{tcolorbox}
\usepackage{enumitem}
\usepackage{array}
\usepackage{fancyhdr}
\usepackage[hidelinks]{hyperref}
\usetikzlibrary{arrows.meta,positioning,calc}
\usepackage[a4paper,margin=2.1cm,headheight=26pt,headsep=10pt]{geometry}
\usepackage{titlesec}

\definecolor{primary}{RGB}{150,98,162}
\definecolor{primarydark}{RGB}{92,52,110}

\tcbset{boxbase/.style={enhanced,breakable,boxrule=0.9pt,arc=2.5pt,left=3mm,right=3mm,top=2.3mm,bottom=2.3mm,fonttitle=\bfseries,coltitle=white,toptitle=0.6mm,bottomtitle=0.6mm}}
\newtcolorbox[auto counter]{exercice}[1][]{boxbase,colback=primary!4,colframe=primary,title={\textsc{Exercice}~\thetcbcounter\ \ifx\relax#1\relax\relax\else\textmd{--- \itshape #1}\fi}}

\newcommand{\dS}{\Delta S}

\pagestyle{fancy}
\fancyhf{}
\fancyhead[L]{\small\textcolor{primarydark}{\textbf{Fiche 8} — Thème 4 : Entropie et deuxième principe}}
\fancyhead[R]{\small\textcolor{primarydark}{\textsc{Moyen}}}
\fancyfoot[C]{\small\thepage}
\renewcommand{\headrulewidth}{0.8pt}
\renewcommand{\headrule}{\hbox to\headwidth{\color{primary}\leaders\hrule height \headrulewidth\hfill}}
\setlength{\parindent}{0pt}\setlength{\parskip}{3pt}

\begin{document}

\begin{center}
\begin{tikzpicture}
\node[rectangle,rounded corners=7pt,fill=primary,text=white,minimum width=16cm,
      minimum height=1.1cm,align=center,inner sep=6pt]
  {{\large\bfseries Thème 4 — Entropie et deuxième principe}\\[1pt]{\small Niveau \textsc{Moyen} — énoncés}};
\end{tikzpicture}
\end{center}

\begin{exercice}[compter les moles de gaz]
Pour chaque réaction, compter les moles de gaz de chaque côté et en déduire le signe de $\dS$.
\begin{enumerate}[label=\alph*),leftmargin=2em,topsep=2pt,itemsep=2pt]
  \item $2\,\ce{SO2}(g) + \ce{O2}(g) \rightarrow 2\,\ce{SO3}(g)$ ;
  \item $\ce{CaCO3}(s) \rightarrow \ce{CaO}(s) + \ce{CO2}(g)$ ;
  \item $\ce{NH3}(g) + \ce{HCl}(g) \rightarrow \ce{NH4Cl}(s)$ ;
  \item $2\,\ce{H2O2}(l) \rightarrow 2\,\ce{H2O}(l) + \ce{O2}(g)$ ;
  \item $\ce{N2}(g) + \ce{O2}(g) \rightarrow 2\,\ce{NO}(g)$.
\end{enumerate}
\end{exercice}

\begin{exercice}[le gaz est le facteur décisif]
Ces réactions contiennent des solides ; c'est le \emph{nombre de moles de gaz} qui tranche. Donner le
signe de $\dS$ (ou $\dS\approx0$) pour chacune.
\begin{enumerate}[label=\alph*),leftmargin=2em,topsep=2pt,itemsep=2pt]
  \item $\ce{C}(s) + \ce{O2}(g) \rightarrow \ce{CO2}(g)$ ;
  \item $2\,\ce{C}(s) + \ce{O2}(g) \rightarrow 2\,\ce{CO}(g)$ ;
  \item $\ce{Fe2O3}(s) + 3\,\ce{CO}(g) \rightarrow 2\,\ce{Fe}(s) + 3\,\ce{CO2}(g)$ ;
  \item $4\,\ce{Fe}(s) + 3\,\ce{O2}(g) \rightarrow 2\,\ce{Fe2O3}(s)$.
\end{enumerate}
\end{exercice}

\begin{exercice}[entropie et réaction réversible]
On considère l'équilibre \;$\ce{N2O4}(g) \rightleftharpoons 2\,\ce{NO2}(g)$.
\begin{enumerate}[label=\alph*),leftmargin=2em,topsep=2pt,itemsep=2pt]
  \item Quel est le signe de $\dS$ dans le \textbf{sens direct} ($\ce{N2O4} \rightarrow 2\,\ce{NO2}$) ?
  \item Quel sens (direct ou inverse) est \emph{favorisé par le désordre} ?
  \item Quel est le signe de $\dS$ dans le sens inverse ?
\end{enumerate}
\end{exercice}

\begin{exercice}[désordre et spontanéité]
À température ambiante (\SI{25}{\celsius}), un glaçon fond spontanément : $\ce{H2O}(s) \rightarrow \ce{H2O}(l)$.
\begin{enumerate}[label=\alph*),leftmargin=2em,topsep=2pt,itemsep=2pt]
  \item Quel est le signe de l'entropie du système $\dS_{\text{système}}$ ?
  \item Cette fusion étant \emph{spontanée} à \SI{25}{\celsius}, que peut-on affirmer sur
        $\dS_{\text{univers}}$ d'après le deuxième principe ?
  \item Le deuxième principe dit qu'un processus spontané « tend à devenir complet ». Qu'observe-t-on
        effectivement pour le glaçon laissé à l'air ?
\end{enumerate}
\end{exercice}

\end{document}
